\documentclass[11pt,a4paper]{article}
\usepackage[utf8]{inputenc}
\usepackage[T1]{fontenc}
\usepackage{hyperref}
\usepackage{graphicx}

\title{Müzikolojinin ve Kognitif Nörobilimin Kesişiminde Sesin Ontolojisi: Estetik, Tarih ve Nörobiyolojik Dinamikler}
\date{}
\maketitle

\tableofcontents
\newpage

\section{Müzikolojinin ve Kognitif Nörobilimin Kesişiminde Sesin Ontolojisi: Estetik, Tarih ve Nörobiyolojik Dinamikler}
\subsection{Müziğin Tanımı, Epistemolojisi ve Kuramsal Temelleri}
Müzik, seslerin belirli ritim, melodi ve armoni kurallarına göre estetik ve duygusal bir biçimde bir araya getirilmesi sanatı olarak tanımlanmaktadır.\textsuperscript{[1]} Basit bir halk şarkısından en karmaşık elektronik kompozisyonlara kadar uzanan bu insan yapımı olgu, tüm kültürel standartlarda kavramsal ve işitsel bir mühendislik ürünü olarak kabul edilir.\textsuperscript{[1]} Antik felsefede Demokritos, müziğin temel bir hayatta kalma ihtiyacından ziyade toplumsal bir refah ve lüksün ("superfluity") sonucu ortaya çıktığını iddia etmiştir.\textsuperscript{[1]} Buna karşın Konfüçyüs, müziği evrensel düzenin ve ahlakın bir aynası olarak konumlandırmış; müziğin hüzün, memnuniyet, sevinç, öfke, dindarlık ve sevgi olmak üzere altı temel duyguyu yansıtarak insan karakterini ifşa ettiğini savunmuştur.\textsuperscript{[1]} Platon ise müziğin kontrolsüz duygusal gücünden endişe duyarak, devletin selameti için tehlikeli müzik modlarının sansürlenmesi gerektiğini ve eğitimin jimnastik ile dengelenmiş doğru müzik tarzlarıyla yapılandırılmasını önermiştir.\textsuperscript{[1]}

Orta Çağ epistemolojisinde müzik, gramer, retorik ve mantıktan oluşan sözel sanatlar (\textit{trivium}) yerine; aritmetik, geometri ve astronomi ile birlikte sayısal ve metafiziksel özelliklere dayanan \textit{quadrivium} kapsamında yüksek bir bilim dalı olarak kabul edilmiştir.\textsuperscript{[4]} Ancak bu yüksek statü zaman zaman sarsılmış, nitekim 1300'lü yıllarda kuramcı Johannes de Grocheo müzik kompozisyonunu dericilik ve ayakkabıcılık gibi sıradan bir el zanaatıyla eş değer tutmuştur.\textsuperscript{[4]}

Batı tonal müziği dikey armoni ve akor yapıları üzerine inşa edilmişken, Güney Asya (Hindistan) klasik müziği tamamen yatay bir melodik karmaşıklığa dayanmaktadır.\textsuperscript{[1]} Hint müziğinde dikey akor yapıları bulunmaz; oktavın bölünmesiyle elde edilen aralıklar Batı müziğine göre çok daha fazla olup melodik karmaşıklık son derece yüksektir.\textsuperscript{[1]} Erken dönemlerde Vedik ilahilerle dini hizmete sunulan bu müzik geleneğinde, anlatıcı (\textit{narrator}) merkezi bir rol oynamakta ve icracının virtüözlüğü enstrümantalistlerle yarışmaktadır.\textsuperscript{[1]}

Yirminci yüzyılda müzikolojide "müzikal ton" kavramı köklü bir kırılma yaşamıştır.\textsuperscript{[1]} Yüzyılın başlarında tonun düzenli titreşimleri ve sabit perdesi onu gürültüden ayırırken; yüzyılın ikinci yarısında John Cage gibi avangart besteciler gürültüyü, rastlantısal sesleri ve sessizliği kompozisyonun meşru unsurları haline getirmiştir.\textsuperscript{[1]} Bu durum, Arthur Schopenhauer ve Friedrich Nietzsche'nin geliştirdiği dinamizm kuramlarıyla da paralellik göstermektedir.\textsuperscript{[5]} Schopenhauer müziği Platonik fikirlerin bir kopyası olarak değil, doğrudan "istencin" (\textit{will}) kendisinin bir yansıması olarak görerek diğer uzamsal sanatlardan üstün tutmuştur.\textsuperscript{[5]} Nietzsche ise müziği Dionysosçu (coşku, esrime, sarhoşluk) ve Apolloncu (form, rasyonellik) dikotomi üzerinden temellendirmiş ve sembol üretiminin insan zihninin otomatik bir faaliyeti olduğunu öne sürmüştür.\textsuperscript{[5]} Bilişsel düzeyde kompozisyon, doğaçlamadan (\textit{improvisation}) tekrarlanabilir bir bütün oluşturması yönüyle ayrılırken; ritmik-melodik en küçük birim olan "motif" (örneğin Beethoven'ın 5. Senfonisi'ndeki kısa-kısa-kısa-uzun yapısı) perdeli ya da perdesiz sesleri zaman içinde düzenleme işlevi görür.\textsuperscript{[4]} Bu kuramsal çerçevede Edmund Gurney'nin \textit{The Power of Sound} (1880) çalışması ve Henri Bergson'un "entelektüel sezgi eylemi" yaklaşımı müziğin bütünsel algısını desteklemektedir.\textsuperscript{[5]}

\subsection{Klasik ve Popüler Müzik Türlerinin Taksonomik Analizi}
Müziğin tarihsel gelişimi, sosyo-kültürel değişimler ve endüstriyel devrimlerle şekillenmiştir.\textsuperscript{[2]} Klasik müzik, kilise ve saray himayesinde katı normlarla gelişirken; popüler müzik endüstrileşme ve kentleşme sonrasında geniş kitlelerin tüketimine sunulan ticari odaklı bir yapı kazanmıştır.\textsuperscript{[2]}

\begin{tabular}{|l|l|l|}
\hline
Klasik Müzik Dönemi & Tarihsel Aralık & Temel Karakteristikler ve Estetik Arayış \\ \hline
\textbf{Medieval (Orta Çağ)} & 500 – 1400 & Gregorian kantoları, tek sesli melodi (unison) ve Katolik Kilisesi odaklı dini kompozisyonlar.\textsuperscript{[7]} \\
\textbf{Renaissance (Rönesans)} & 1400 – 1600 & Polifoni (çok seslilik) kullanımı ve bağımsız melodilerin eş zamanlı bir uyum içinde bir araya getirilmesi.\textsuperscript{[7]} \\
\textbf{Baroque (Barok)} & 1600 – 1750 & Süslü melodiler, karmaşık armoniler ve kontrast yaratan yoğun enstrümantasyon yapıları.\textsuperscript{[7]} \\
\textbf{Galant} & 1720 – 1770 & Barok dönemin yoğun karmaşıklığına bir tepki olarak doğan sade, zarif ve melodi odaklı geçiş tarzı.\textsuperscript{[7]} \\
\textbf{Classical (Klasik)} & 1750 – 1820 & Formda denge, netlik, simetri ve yapısal sadelik arayışı.\textsuperscript{[7]} \\
\textbf{Romantic (Romantik)} & 1780 – 1910 & Yoğun duygu, bireysel ifade, dramatik anlatım ve kural dışı form arayışları.\textsuperscript{[7]} \\
\textbf{Modernism / High Modernism} & 1890 – Günümüz & Geleneksel tonalite ve formların reddi, atonalite, serializm ve radikal deneysel teknikler.\textsuperscript{[7]} \\
\textbf{Postmodern \& Experimental} & 1930 / 1950 – Günümüz & Eski ve yeninin kolajı, kuralların esnetilmesi, şans faktörü ve gürültü ile sessizliğin entegrasyonu.\textsuperscript{[1]} \\
\hline
\end{tabular}
Popüler müziğin evrimi ise geleneksel sözlü aktarıma dayanan halk müziğinin gerilemesiyle hız kazanmıştır.\textsuperscript{[2]} 1890'larda New York'ta kurulan Tin Pan Alley, tarihin ilk organize şarkı yayıncılık endüstrisini oluşturmuş ve Avrupa operetiyle birleşerek Broadway müzikallerinin (Jerome Kern, George Gershwin, Irving Berlin, Cole Porter, Richard Rodgers, Oscar Hammerstein II) doğuşunu sağlamıştır.\textsuperscript{[9]} Teknolojik kırılmalar, özellikle fonograf plaklarının evlere girmesi ve mikrofonun icadı, vokalistlerin geniş konser salonlarında seslerini duyurmak için bağırmak zorunda kalmasını engellemiş, böylece "crooner" adı verilen daha samimi ve içsel vokal tekniklerinin (Frank Sinatra, Bing Crosby) gelişmesine olanak tanımıştır.\textsuperscript{[9]}

1940'larda güneyli siyah nüfusun kuzey sanayi kentlerine göç etmesiyle (Büyük Göç), geleneksel kırsal blues tarzı (Georgia/Carolinas, Texas ve Mississippi Delta stilleri) şehre uyarlanarak Chicago blues adı verilen elektrikli, sert ve tam ritim grubu destekli bir form kazanmıştır.\textsuperscript{[9]} Bu elektro-gitar tabanlı altyapı, jazz tınıları ile birleşerek ritim ve blues (R\&B) tarzını, o da country müzikle melezlenerek 1950'lerde gençlik isyanının sembolü haline gelen Rock and Roll'u (Elvis Presley, Chuck Berry) doğurmuştur.\textsuperscript{[2]} 1960'larda Beatles ve Rolling Stones ile İngiliz İstilası (\textit{British Invasion}) yaşanmış, Bob Dylan ile Folk Rock entelektüel derinlik kazanmıştır.\textsuperscript{[9]}

1970'ler ve 1980'lerde rock müzik; disco, punk ve heavy metal gibi alt türlere ayrışırken, Michael Jackson'ın 1982 tarihli \textit{Thriller} albümü pop, rock ve R\&B'yi sentezleyerek müzik videoları üzerinden popüler müziği çok boyutlu bir multimedya çağına taşımıştır.\textsuperscript{[2]} Eş zamanlı olarak Bronx'ta doğan hip-hop ve rap, DJ Kool Herc'ün pikaplar üzerinde "breakbeat" oluşturmasıyla ritmik bir devrim yaratmıştır.\textsuperscript{[9]} Country müzik ise Appalachians göçmenlerinin İngiliz-İrlanda baladlarından türemiş, radyo yayınları (\textit{Grand Ole Opry}) ve Nashville Sound'u (Chet Atkins) ile ticarileşmiş, Outlaw Country (Willie Nelson) ve 1990'larda Garth Brooks'un stadyum konserleri seviyesindeki şovlarıyla küreselleşmiştir.\textsuperscript{[9]}

\begin{tabular}{|l|l|l|l|}
\hline
Popüler Müzik Türü & Kökeni ve Tarihsel Kırılma Noktası & Temel Enstrümantasyon ve Yapısal Form & Sosyo-Kültürel ve Teknolojik Dinamikler \\ \hline
\textbf{Jazz} & New Orleans (19. yy sonu - 20. yy başı) \textsuperscript{[9]} & Trompet, trombon, saksofon, piyano, kontrbas; senkoplu ritimler ve doğaçlama.\textsuperscript{[9]} & Afrika ritimleri ile Avrupa armonilerinin sentezi; swing büyük orkestraları ve radyo dönemi.\textsuperscript{[9]} \\
\textbf{Blues} & Güney ABD Plantasyonları (Amerikan İç Savaşı sonrası) \textsuperscript{[9]} & Akustik/elektrik gitar, mızıka; çağrı-cevap yapısı ve mavi notalar.\textsuperscript{[9]} & Siyah Amerikalıların tarlalardaki iş şarkıları ve feryatları; "Race Records" plak sektörü.\textsuperscript{[9]} \\
\textbf{Rock and Roll / Rock} & ABD (1950'lerin ortası) \textsuperscript{[9]} & Elektro-gitar, bas gitar, davul; sert arka ritim (backbeat) ve yüksek ses gücü.\textsuperscript{[9]} & R\&B ve country sentezi; gençlik isyanı, amplifikatör teknolojisi ve gençlerin harcama gücü.\textsuperscript{[9]} \\
\textbf{Country} & Güney ve Batı ABD Kırsalı (20. yy başı) \textsuperscript{[9]} & Keman, banjo, gitar, mandolin, çelik gitar; basit akorlar ve yodel tekniği.\textsuperscript{[9]} & İngiliz, İskoç ve İrlandalı yerleşimcilerin halk baladları; göç dalgaları ve radyo yayılımı.\textsuperscript{[9]} \\
\textbf{Hip-Hop / Rap} & Bronx, New York (1970'lerin sonu) \textsuperscript{[9]} & Turntable (pikap), sampler, drum machine; ritmik konuşma (rapping) ve breakbeat.\textsuperscript{[9]} & DJ Kool Herc'ün Jamaika ses sistemlerini getirmesi; azınlık topluluklarının sosyo-politik sesi.\textsuperscript{[9]} \\
\hline
\end{tabular}
\subsection{Ses Algısının Fizyolojik ve Nöroanatomik Kodlaması}
Ses dalgalarının algılanması, dış kulaktan başlayıp serebral korteksin en derin bölgelerine kadar uzanan karmaşık bir nörolojik koordinasyon gerektirir.\textsuperscript{[14]} Kulak zarına çarpan mekanik dalgalar, orta kulaktaki kemikçikler (özellikle stapes) vasıtasıyla spiral yapılı kohleaya (salyangoz) iletilir.\textsuperscript{[15]} Kohleanın içindeki sıvı dalgalanmaları, yaklaşık 10.000 ila 15.000 adet tüy hücresinde (cilia) mekanik-elektriksel dönüşüm (transdüksiyon) gerçekleştirerek kimyasal nörotransmitterlerin salınmasını sağlar ve işitsel siniri uyarır.\textsuperscript{[15]} Buradan çıkan elektrik akımları beyin sapı üzerinden temporal lobdaki birincil işitsel kortekse ulaştırılır.\textsuperscript{[14]}

\begin{tabular}{|l|l|l|}
\hline
Beyin Bölgesi / Yapısı & Spesifik Nöral Görevi & Fonksiyonel Analiz ve İşlevi \\ \hline
\textbf{Sağ Temporal Lob} & Perde, melodi ve armoni çözümlemesi.\textsuperscript{[15]} & Notaların perde dizilimlerini, akorları ve dikey harmonik yapıları algılar.\textsuperscript{[15]} \\
\textbf{Temporal Gyrus (Yakın Alanlar)} & Timbre (tını) ayrımı.\textsuperscript{[15]} & Aynı notayı çalan farklı enstrümanların (örneğin keman ile piyano) ses kalitesini ayırt eder.\textsuperscript{[15]} \\
\textbf{Serebellum (Beyincik)} & Ritim, zamanlama ve motor koordinasyon.\textsuperscript{[14]} & Müziğin zamansal yapısını çözer, motor sistemle senkronizasyon kurarak fiziksel hareketi düzenler.\textsuperscript{[14]} \\
\textbf{Amigdala ve Hipokampus} & Duygusal hafıza ve otobiyografik bellek.\textsuperscript{[14]} & Müziğe bağlı anıları çağırır ve bu anıların tetiklediği duygusal reaksiyonları yönetir.\textsuperscript{[14]} \\
\textbf{Mezolimbik Ödül Sistemi} & Dopaminerjik aktivasyon ve haz duyumu.\textsuperscript{[14]} & Güçlü müzikal uyarılma anlarında ("spine-tingling" ürperme hissi) dopamin salgılatır.\textsuperscript{[15]} \\
\textbf{Motor Sistem} & Ritim tahmini ve kas senkronizasyonu.\textsuperscript{[17]} & Bireg henüz fiziksel olarak ritim tutmaya başlamadan önce vuruşu (beat) tahmin ederek motor sistemi aktive eder.\textsuperscript{[17]} \\
\textbf{Orbitofrontal Korteks (OFC)} & Beklenti, tahmin ve gerilim-çözülme analizi.\textsuperscript{[16]} & Western tonal müziğindeki gerilim (tension) ve çözülme (resolution) kalıplarını işler; OKB hastalarında hiperaktivite gösterir.\textsuperscript{[16]} \\
\hline
\end{tabular}
İşitsel uyaranların beyinde yarattığı etki sadece bu bölgelerle sınırlı kalmayıp otonom sinir sistemini (ANS) de doğrudan manipüle etmektedir.\textsuperscript{[16]} Müziğin pozitif veya negatif hissettirme derecesi (valans) ve ritmik uyarıcılığı kalp atış hızı ve solunum sıklığı gibi istemsiz süreçleri değiştirir.\textsuperscript{[16]} Örneğin sinematik yapıtlarda (Interstellar, Schindler's List, Lord of the Rings, Mulan veya Jaws gibi filmlerde) gerilim-çözülme kalıplarının kullanılması, izleyicinin orbitofrontal korteksini uyararak beklenti ve tahmin mekanizmalarını tetikler ve kalp atış ritmini doğrudan hızlandırır.\textsuperscript{[16]}

\subsection{Klinik Rehabilitasyon, Bilişsel Kapasite ve Metodolojik Sınırlar}
Müzikal uyaranlar beynin geniş ağlarını aktive ettiğinden, nöroplastisiteyi tetikleme potansiyelleri son derece yüksektir.\textsuperscript{[14]} Örneğin uzun süreli enstrüman çalanlarda (keman çalanların yay parmaklarıyla ilişkili beyin alanlarındaki sinaptik yoğunluk gibi), ilgili frekanslara duyarlı kortikal bölgelerin sinaptik yapısının zenginleştiği gözlemlenmiştir.\textsuperscript{[14]} Kognitif gelişim alanında, UC Irvine araştırmacılarının ortaya koyduğu "Mozart Etkisi", Mozart'ın Re Majör İki Piyano Sonatı'nı (K. 448) 10 dakika dinlemenin, uzamsal-zamansal akıl yürütme skorlarını geçici olarak (yaklaşık 15 dakika boyunca) 8 ila 9 IQ puanı kadar artırdığını göstermiştir.\textsuperscript{[15]} Bu sonatın ayrıca epilepsi hastalarında nöbet sıklığını azalttığı saptanmıştır.\textsuperscript{[16]}

Beyin hasarları, müzikal algıda "sensory amusia" (müziği algılama veya tepki verme yetisinin kaybı) veya "musical hallucinosis" (dışarıda bir müzik çalmadığı halde sürekli müzik duyma illüzyonu) gibi spesifik patolojilere yol açabilir.\textsuperscript{[14]} Nörolog Oliver Sacks, \textit{Musicophilia} adlı çalışmasında bu tip amuzi ve sıra dışı müzikal algı bozukluklarını geniş bir vaka analiziyle incelemiştir.\textsuperscript{[15]} Bu hasarların klinik rehabilitasyonunda müziğin iyileştirici gücü somut verilerle desteklenmektedir:

\begin{tabular}{|l|l|l|}
\hline
Hedef Klinik Durum & Uygulanan Müzikal Müdahale & Gözlemlenen Fizyolojik ve Klinik Sonuçlar \\ \hline
\textbf{Sol Hemisfer Hasarı / Afazi} & Melodik İntonasyon Terapisi (konuşmayı şarkı formuna dökme).\textsuperscript{[18]} & Hasarsız olan sağ beyni uyararak konuşma yetisinin geri kazanılması (örneğin Gabrielle Giffords'un iyileşme süreci).\textsuperscript{[18]} \\
\textbf{Alzheimer ve Demans} & Kişiselleştirilmiş nostaljik çalma listelerinin dinletilmesi.\textsuperscript{[14]} & Uzun süreli bellekte kodlanmış "Happy Birthday" gibi şarkılarla bellek kapılarının açılması, geçici bilişsel uyanış ve sosyal iletişim artışı.\textsuperscript{[14]} \\
\textbf{Parkinson Hastalığı} & Ritmik adım atma ve hareket-müzik senkronizasyonu.\textsuperscript{[14]} & Serebellar ve motor devrelerin uyarılmasıyla yürüme donmalarının önüne geçilmesi ve fiziksel mobilitenin artması.\textsuperscript{[14]} \\
\textbf{Düşme Riski Olan Yaşlılar} & Altı ay boyunca müzik eşliğinde yürüyüş ve ritim eğitimi.\textsuperscript{[15]} & Motor koordinasyon ve dengenin iyileşmesi; kontrol grubuna kıyasla düşme vakalarında \%54 oranında düşüş.\textsuperscript{[15]} \\
\textbf{Cerrahi Hastaları (Pre/Post-op)} & Yavaş tempolu enstrümantal müzik dinletimi.\textsuperscript{[15]} & Kan basıncında (sistolik , diyastolik ) düşüş, adrenaline ve IL-6 düzeylerinde azalma, sedatif ilaç gereksiniminde düşüş.\textsuperscript{[15]} \\
\hline
\end{tabular}
Ancak bu iddiaların ve saha araştırmalarının metodolojik sınırları da göz önünde bulundurulmalıdır.\textsuperscript{[17]} Örneğin, AARP tarafından gerçekleştirilen ve erken yaşta müzikle tanışmayan ancak yetişkinlikte müzikle ilgilenen bireylerin daha yüksek bilişsel refah skorlarına sahip olduğunu öne süren geniş çaplı anket çalışması (3.185 ABD'li yetişkin) önemli metodolojik açmazlar barındırmaktadır.\textsuperscript{[17]} Çalışmanın tamamen öznel anket beyanlarına dayanması, MRI gibi objektif nörogörüntüleme veya standardize kognitif testler içermemesi güvenirliği sınırlamaktadır.\textsuperscript{[17]} En önemlisi, nedensellik ilişkisinin net olmamasıdır; nitekim çocukken müzik eğitimi alan bireylerin yetişkinlikteki bilişsel üstünlüğü müzikten ziyade, bu eğitimi almalarını sağlayan yüksek sosyo-ekonomik aile yapısı ve nitelikli eğitim imkanlarından kaynaklanıyor olabilir.\textsuperscript{[17]} Beyin patikalarının kullanılmadığında köreldiği ("use it or lose it") göz önüne alındığında, müziğin bu geniş nöral aktivasyon gücü kognitif rezervi korumada kritik bir rol oynamaya devam etmektedir.\textsuperscript{[17]}

\subsection{Sonuç}
Müzik, insanlığın sadece kültürel bir lüksü ya da estetik bir uğraşı değil, aynı zamanda sinir sisteminin gelişimini, organizasyonunu ve onarımını doğrudan etkileyen biyolojik bir katalizördür.\textsuperscript{[1]} Antik dönemden bu yana felsefi tartışmalara konu olan müziğin insan davranışları üzerindeki dönüştürücü gücü, günümüzde modern nörogörüntüleme teknikleri ve klinik çalışmalarla somut birer veri haline gelmiştir.\textsuperscript{[14]} Müziğin beyindeki işitsel, motor, bilişsel ve duygusal merkezleri bütünsel olarak senkronize etme yeteneği, onu tıp, eğitim ve sosyal alanlarda benzersiz bir terapötik araç kılmaktadır.\textsuperscript{[14]} Özellikle yaşlanan nüfusta demans yönetiminde, inme sonrası nörolojik rehabilitasyonda ve cerrahi stres kontrolünde müzik tabanlı müdahalelerin standart tedavi protokollerine entegre edilmesi, hem hasta yaşam kalitesini artıracak hem de farmakolojik yükü hafifletecektir.\textsuperscript{[14]}


\bigskip
\noindent\rule{\textwidth}{0.4pt}
\subsection*{Sources}
\noindent [1] Music | Art Form, Styles, Rhythm, \& History | Britannica --- \url{https://www.britannica.com/art/music}\\
\noindent [2] Popular music | Description, History, \& Facts | Britannica --- \url{https://www.britannica.com/art/popular-music}\\
\noindent [4] Musical composition | Definition, History, Structure, Types, \& Facts | Britannica --- \url{https://www.britannica.com/art/musical-composition}\\
\noindent [5] Music - Theories of musical meaning since the 19th century | Britannica --- \url{https://www.britannica.com/art/music/Theories-of-musical-meaning-since-the-19th-century}\\
\noindent [7] Music History Timeline: Evolution of Music Overview --- \url{https://www.lamaisonschoolofmusic.com/resources/music-history-timeline}\\
\noindent [9] Popular art - Music Genres, Culture, Trends | Britannica --- \url{https://www.britannica.com/art/popular-art/Popular-music}\\
\noindent [14] Music and the Brain | Harvard Medical School --- \url{https://hms.harvard.edu/news-events/publications-archive/brain/music-brain}\\
\noindent [15] Music and health - Harvard Health --- \url{https://www.health.harvard.edu/healthy-aging-and-longevity/music-and-health}\\
\noindent [16] How Music Resonates in the Brain | Harvard Medicine Magazine --- \url{https://magazine.hms.harvard.edu/articles/how-music-resonates-brain}\\
\noindent [17] Why is music good for the brain? - Harvard Health --- \url{https://www.health.harvard.edu/blog/why-is-music-good-for-the-brain-2020100721062}\\
\noindent [18] Music can boost memory and mood - Harvard Health --- \url{https://www.health.harvard.edu/mind-and-mood/music-can-boost-memory-and-mood}\\

\end{document}
